\documentclass[11pt]{article}
\usepackage{setspace}
\usepackage{geometry}
\geometry{margin=1in}
\setstretch{1.15}

\begin{document}

\noindent
\textbf{Zhongli Wang} \\
+49 160 9541 6015 \\
nkwangzhongli@gmail.com \\
https://sites.google.com/view/zhongli-wang \\
Bielefeld, Germany \\

\vspace{0.5cm}

\noindent
\textbf{Zeppelin Universität} \\
Personalmanagement -- Ute Aisenpreis \\
Friedrichshafen, Germany \\
% (Parenthood leave replacement until end of September 2027, 75\%) \\

\vspace{0.5cm}

\noindent
\textbf{Coordination of Academic Programs (m/w/d) -- 75\% (bis 09/2027)} \\

\vspace{0.3cm}

\noindent Dear Ms.\ Aisenpreis,\\[0.2cm]

I am writing to apply for the position \textit{Coordination of Academic Programs (m/w/d)} (75\%) at Zeppelin Universität. 
I am currently completing my Ph.D.\ in Economics and have gained hands-on experience working at the interface of teaching, administration, and data-driven process improvement. 
I am attracted to this role because it combines structured coordination work with meaningful impact on students and teaching staff, and I would be glad to support the ZU team in stabilising and further developing campus management processes during the parenthood leave replacement period.

During my Master’s studies at the University of Zurich, I worked for two years as a teaching assistant at the Chair of Quantitative Business Administration and contributed to the SYLVA project—an automated scoring and feedback system used in large courses. In this role, I supported the preparation and evaluation of exams by helping translate course requirements into a reliable assessment workflow, improving consistency, scalability, and transparency in grading and feedback delivery. This experience strengthened my appreciation for well-designed academic support systems and gave me hands-on exposure to the practical interface between teaching operations, digital tools, and student-facing processes—experience that maps directly onto the program coordination and campus-management responsibilities of this position.

During my PhD, I have led multiple end-to-end projects that required exactly the kind of structured coordination, precision, and systems thinking that this position demands. I designed and maintained reproducible research workflows that integrated data collection, version control, documentation, and automated reporting; coordinated timelines and deliverables across international partners within the Marie Skłodowska-Curie EPOC network; and translated complex requirements into implementable structures (e.g., mapping program rules into clear course structures, checks, and templates for consistent execution). This background is not separate from program management—it is strong evidence of my ability to handle complex rule-based systems reliably, communicate across stakeholder groups, and ensure high-quality outcomes under deadlines. Applying for the Coordination of Academic Programs role therefore reflects a well-matched skill set and a professional interest in supporting high-quality teaching operations with analytical rigor, rather than a departure from my academic profile.

My background fits the core requirements of the position in three ways. 
First, I am highly comfortable with complex rule-based systems and structured workflows. 
In my academic work, I routinely translate formal requirements into implementable structures---from course and program logic to data pipelines and documentation standards. 
This is directly relevant for modelling and maintaining study and examination regulations in a Campus Management System, checking feasibility across programmes, and ensuring consistent ``W\"ahlbarkeit'' and correct implementation.

Second, I bring strong technical affinity and a practical mindset for systems work. 
I work confidently with digital tools, databases and structured data, and I learn new software quickly. 
In a CMS setting, I would be ready to act as a reliable key user: documenting recurring issues, proposing clear fixes, supporting testing cycles, and building user-friendly guides for colleagues. 
I would also be happy to contribute to CMS trainings for staff and to the maintenance of relevant intranet pages, with an emphasis on clarity and usability.

Third, I have substantial experience in coordination and communication across diverse stakeholders. 
Through teaching-related work, workshops, and research team coordination, I have learned to align timelines, communicate decisions clearly, and handle many small dependencies without losing oversight. 
In this position, I would be glad to support course-choice preparation and execution end-to-end, coordinate interfaces with the Studies and Examination Center, and ensure smooth communication within and across departments.

I am fluent in English and have advanced German proficiency (C1). 
I also value a collaborative working culture and would be pleased to contribute beyond the minimum: proactively taking ownership of process documentation, supporting the core A5 (CMS) team where needed, and helping ensure that programme structures and course-choice processes remain robust, transparent, and student-friendly.

Thank you for considering my application. 
I would welcome the opportunity to discuss how I can support Zeppelin Universität with reliable coordination, system-focused implementation, and clear communication.\\[0.4cm]

\noindent Sincerely,\\[1.2cm]
\noindent Zhongli Wang

\end{document}